\documentclass[11pt,a4paper]{article}
\usepackage{turkos}

\begin{document}
\phaseheader{0}{Foundations and Toolchain}{Set up the workshop: a cross compiler, an emulator, a debugger, and the C habits a kernel needs}{0}{5}

\ataglance{2--3 weeks}{1}{Basic C, a Fedora Linux machine, curiosity}{A working \texttt{i686-elf} cross toolchain, QEMU and GDB,
a Git repository with the Turk-OS skeleton, and a small C library (strings, bitmaps, lists, ring buffer) with unit tests.
You will reuse that library inside the kernel.}

\section{Why this phase}

An operating system kernel is a C program with nothing underneath it: no standard library, no \code{printf}, no
\code{malloc}, no operating system to catch its mistakes. Before you can write one, you need three things. The first
is a \textbf{toolchain} that produces code for a bare x86 machine instead of for Linux. The second is a
\textbf{laboratory}, an emulator you can reboot a hundred times an hour with a debugger attached. The third is
\textbf{C fluency at the level of bits and addresses}: pointers to fixed memory locations, bit masks, struct layout,
and the build pipeline from \code{.c} file to executable.

This phase builds all three and gives you the vocabulary of the field. It is the only phase with no kernel code, and
the temptation is to rush it. Don't. Almost every mysterious failure in Phases 1--3 comes from a toolchain
misconfiguration or a C misunderstanding that could have been fixed here in ten minutes.

\section{Concepts you need}

\subsection{What an operating system is}

\OSTEP describes the OS through three ideas that also organise the whole book and this plan. It
\textbf{virtualizes} the hardware, so each program believes it has its own CPU and its own memory (Phases 5--11). It
manages \textbf{concurrency}, so many activities can safely share the machine (Phases 8--9). It provides
\textbf{persistence}, so data survives power loss (Phases 12--13). \MOS gives two complementary views of the same
thing: the OS as an \emph{extended machine} that hides ugly hardware behind clean abstractions, and as a
\emph{resource manager} that shares the CPU, memory and devices fairly and safely.

\begin{conceptbox}{Kernel mode and user mode}
The CPU has at least two privilege levels. Code in \textbf{kernel mode} (ring 0 on x86) can do anything: touch any
memory, talk to devices, disable interrupts. Code in \textbf{user mode} (ring 3) can't, and if it tries, the CPU
stops it and hands control to the kernel. This hardware-enforced split is what lets an OS protect itself and its
programs from each other (\OSC\,\S1.4). For the first nine phases everything you write runs in ring 0. Ring 3 arrives
in Phase 10.
\end{conceptbox}

\subsection{The machine you are going to program}

Read \MOS\,\S1.3 with this picture in mind. The \textbf{CPU} fetches instructions from memory and executes them,
keeping its working values in a handful of \textbf{registers}. \textbf{Memory} is one big array of bytes, indexed by
address. \textbf{Devices} (screen, keyboard, disk, timer) are controlled by reading and writing their registers,
either through a separate I/O address space with special \code{in}/\code{out} instructions, or by mapping them into
ordinary memory addresses. When a device needs attention it raises an \textbf{interrupt}, which makes the CPU pause
what it is doing and jump to a handler the OS registered earlier. Nearly everything in Phases 2--4 is a concrete case
of these four sentences.

\subsection{C the way a kernel uses it}

A kernel is compiled \textbf{freestanding}: there is no C library, but the compiler still provides the headers that
only define types and macros: \texttt{stdint.h}, \texttt{stddef.h}, \texttt{stdbool.h}, \texttt{stdarg.h},
\texttt{limits.h} and a few others. Everything else you write yourself. The listing shows the idioms you will use every
day. Make sure you can explain each line before moving on.

\begin{lst}{c}{idioms.c (study this, you will write code like it constantly)}
#include <stdint.h>
#include <stddef.h>

/* 1. A pointer to a fixed physical address: the VGA text buffer.            */
/*    volatile = "every read/write really happens"; the compiler may not     */
/*    cache, merge or delete these accesses.                                  */
volatile uint16_t *const vga = (volatile uint16_t *)0xB8000;

/* 2. Bit manipulation: set, clear, test bit n of a flags word.              */
#define BIT(n)            (1u << (n))
#define SET(x, n)         ((x) |=  BIT(n))
#define CLEAR(x, n)       ((x) &= ~BIT(n))
#define TEST(x, n)        (((x) & BIT(n)) != 0)

/* 3. Alignment: round an address up/down to a power-of-two boundary.        */
#define ALIGN_DOWN(x, a)  ((x) & ~((uintptr_t)(a) - 1))
#define ALIGN_UP(x, a)    ALIGN_DOWN((x) + (a) - 1, (a))

/* 4. A struct whose memory layout is dictated by hardware: no padding.      */
struct __attribute__((packed)) gdt_ptr {
    uint16_t limit;
    uint32_t base;
};                       /* sizeof == 6, not 8 */

/* 5. Function pointers: tables of handlers, e.g. one per interrupt.         */
typedef void (*handler_t)(void);
static handler_t handlers[256];
\end{lst}

\subsection{From source file to bootable image}

You will spend a lot of time looking at what the toolchain produced, so learn the pipeline now
(Figure~\ref{fig:pipeline}). The compiler turns C into assembly and the assembler turns that into an
\textbf{object file} containing \textbf{sections}: \code{.text} (code), \code{.rodata} (constants), \code{.data}
(initialised globals) and \code{.bss} (zero-initialised globals, which take no space in the file). The
\textbf{linker} combines object files and places each section at the address a \textbf{linker script} tells it to.
The result is an \textbf{ELF} executable that GRUB or QEMU can load. \OSC\,\S2.5 covers linkers and loaders.

\begin{figure}[htbp]
\centering
\begin{tikzpicture}[node distance=0.55cm and 0.7cm, font=\sffamily\small]
  \node[tkbox] (c) {\texttt{kmain.c}};
  \node[tkbox, right=of c] (cc) {\texttt{i686-elf-gcc}\\\scriptsize preprocess + compile};
  \node[tkbox, right=of cc] (s) {\texttt{kmain.o}};
  \node[tkbox, below=of c] (asm) {\texttt{loader.s}};
  \node[tkbox, right=of asm] (nasm) {\texttt{nasm -f elf32}};
  \node[tkbox, right=of nasm] (o2) {\texttt{loader.o}};
  \node[tkhi, right=1.0cm of $(s)!0.5!(o2)$] (ld) {\texttt{i686-elf-ld}\\\scriptsize + \texttt{link.ld}};
  \node[tkbox, right=of ld] (elf) {\texttt{kernel.elf}};
  \node[tkok, below=0.6cm of elf] (q) {QEMU / GRUB};
  \draw[tkarrow] (c) -- (cc); \draw[tkarrow] (cc) -- (s);
  \draw[tkarrow] (asm) -- (nasm); \draw[tkarrow] (nasm) -- (o2);
  \draw[tkarrow] (s) -- (ld); \draw[tkarrow] (o2) -- (ld);
  \draw[tkarrow] (ld) -- (elf); \draw[tkarrow] (elf) -- (q);
\end{tikzpicture}
\caption{The build pipeline you set up in this phase and use for the rest of the project.}
\label{fig:pipeline}
\end{figure}

\subsection{Numbers you will see every day}

Kernel code talks in hexadecimal and powers of two. Memorise this table; you will recognise these values in register
dumps long before you understand everything around them.

\begin{center}
\small\renewcommand{\arraystretch}{1.25}
\begin{tabular}{@{}l l l@{}}
\toprule
\tkth{Value} & \tkth{Hex} & \tkth{Where you meet it} \\
\midrule
512 bytes & \texttt{0x200} & one disk sector, the boot sector (Phase 1) \\
4 KiB & \texttt{0x1000} & one page or page frame (Phases 5--6) \\
1 MiB & \texttt{0x100000} & where the kernel is loaded (Phase 1) \\
4 MiB & \texttt{0x400000} & memory covered by one page table (Phase 6) \\
3 GiB & \texttt{0xC0000000} & start of the higher-half kernel (Phase 6) \\
4 GiB & \texttt{0x100000000} & the whole 32-bit address space \\
\bottomrule
\end{tabular}
\end{center}

\section{Reading plan}

Read the \must items before starting the tasks. The \should items make the tasks easier and are worth the time
this early. Your reading order: \OSTEP first, then \MOS\,\S1.3, then \LOB.

\begin{readingplan}
\must & \OSTEP & Ch.~1 \emph{A Dialogue on the Book}; Ch.~2 \emph{Introduction to Operating Systems} & 1--22 & The three big ideas; the map for the whole plan. \\
\must & \MOS & \S1.3 \emph{Computer Hardware Review} (processors, memory, I/O devices, buses, booting) & 20--35 & The hardware you will program directly. \\
\must & \LOB & Ch.~1 \emph{Introduction}; Ch.~2 up to \emph{Booting} (Tools, Quick Setup, Build System, Virtual Machine) & 7--12 & How the practical guide works and what it expects. \\
\should & \MOS & \S1.1 \emph{What Is an OS?}; \S1.5 \emph{OS Concepts} & 3--6, 38--50 & Vocabulary: processes, address spaces, files, shell. \\
\should & \MOS & \S1.8 \emph{The World According to C} & 73--77 & Header files, large projects, the C runtime model. \\
\should & \OSC & Ch.~1 \S1.1--1.4 (organization, interrupts, storage, dual-mode) & 4--27 & A second, more detailed view of the same ground. \\
\could & \OSC & \S1.9 \emph{Kernel Data Structures}; \S2.5 \emph{Linkers and Loaders} & 36--40, 75--77 & Background for the warm-up library and the build pipeline. \\
\could & \OSDI & \S1.1--1.2 (what an OS is, history, history of MINIX 3) & 4--19 & Where MINIX came from and why it matters for teaching. \\
\end{readingplan}

\section{Build it}

\task{Install the host tools}
Everything below is packaged in Fedora. \code{grub2-tools-extra} provides \code{grub2-mkrescue} (Phase 1) and
\code{mtools} is used by it. The second line installs what GCC needs to build from source.

\begin{lst}{sh}{terminal (host)}
sudo dnf install gcc make git nasm qemu-system-x86 gdb xorriso mtools \
                 grub2-tools-extra grub2-pc-modules
sudo dnf install bison flex gmp-devel mpfr-devel libmpc-devel texinfo
\end{lst}
\verify{\code{nasm -v}, \code{qemu-system-i386 --version} and \code{gdb --version} all print a version.}

\task{Make your editor's terminal see those tools}
Your VS Code is the Flatpak build, so its integrated terminal runs inside a sandbox where host programs like
\code{nasm} and \code{qemu-system-i386} are not visible. Either work in a normal terminal window, or add a terminal
profile to VS Code's \code{settings.json} that starts a shell on the host:

\begin{lst}{plain}{settings.json (VS Code)}
"terminal.integrated.profiles.linux": {
  "host-bash": { "path": "/usr/bin/flatpak-spawn", "args": ["--host", "--env=TERM=xterm-256color", "bash"] }
},
"terminal.integrated.defaultProfile.linux": "host-bash"
\end{lst}
\verify{In a new VS Code terminal, \code{which nasm} prints \code{/usr/bin/nasm}.}

\task{Build the i686-elf cross compiler}
The Fedora \code{gcc} produces programs for Linux. It assumes things that are false inside your kernel: that a C
library exists, that the stack protector runtime is available, that code should be position-independent. A
\textbf{cross compiler} targeting \code{i686-elf} assumes nothing. Follow the OSDev Wiki page \emph{GCC
Cross-Compiler}. Download recent binutils and GCC sources from \url{https://ftp.gnu.org/gnu/}, then build each one in
a separate build directory, never inside its source tree:

\begin{lst}{sh}{build-cross.sh (outline; replace X.Y with the versions you downloaded)}
export PREFIX="$HOME/opt/cross"
export TARGET=i686-elf
export PATH="$PREFIX/bin:$PATH"

mkdir -p ~/src/build-binutils && cd ~/src/build-binutils
../binutils-X.Y/configure --target=$TARGET --prefix="$PREFIX" \
    --with-sysroot --disable-nls --disable-werror
make -j"$(nproc)" && make install

mkdir -p ~/src/build-gcc && cd ~/src/build-gcc
../gcc-X.Y.Z/configure --target=$TARGET --prefix="$PREFIX" \
    --disable-nls --enable-languages=c --without-headers
make -j"$(nproc)" all-gcc all-target-libgcc
make install-gcc install-target-libgcc
\end{lst}

Add \code{export PATH="$HOME/opt/cross/bin:$PATH"} to \code{~/.bashrc}. The GCC build takes 20--40 minutes.

\begin{tipbox}[If the cross build fails]
You can start Phase 1 with the host compiler while you fix it, as \LOB does. Use these flags and link with
\code{ld -m elf_i386}. Switch to the cross compiler before Phase 3, because the host compiler's hidden assumptions
cause bugs that are very hard to recognise.
\begin{lst}{plain}{fallback CFLAGS for the host gcc}
-m32 -ffreestanding -fno-pie -fno-stack-protector -fno-builtin -nostdlib -fcf-protection=none
\end{lst}
\end{tipbox}
\verify{\code{i686-elf-gcc --version} and \code{i686-elf-ld --version} work in a fresh terminal.}

\task{Create the repository skeleton}
Create the project layout, initialise Git, and ignore build outputs. Commit early and often from now on: small
commits let you bisect your way back to the last working kernel when something breaks.

\begin{lst}{sh}{terminal}
cd Turk-OS                     # your project folder
mkdir -p boot kernel lib include tests tools user
git init
printf 'build/\n*.o\n*.elf\n*.bin\n*.iso\n*.img\n' > .gitignore
git add . && git commit -m "Turk-OS: empty skeleton and study plan"
\end{lst}
\verify{\code{git log} shows your first commit.}

\task{Write the C warm-up library, and test it on Linux}
Write these functions from scratch, without the C library, in \code{lib/}. They are small, and every one of them
reappears in the kernel. Compile them with the \emph{host} compiler for now and test them with \code{assert}, using the
address and undefined-behaviour sanitizers to catch memory mistakes while you still have an OS to catch them.

\begin{center}
\small\renewcommand{\arraystretch}{1.25}
\begin{tabularx}{\linewidth}{@{}l X l@{}}
\toprule
\tkth{File} & \tkth{Functions} & \tkth{Used again in} \\
\midrule
\code{string.c} & \code{memset}, \code{memcpy}, \code{memmove}, \code{memcmp}, \code{strlen}, \code{strcmp}, \code{strncpy} & Phase 2 (GCC may call these itself) \\
\code{convert.c} & \code{utoa(value, buf, base)} for bases 10 and 16, \code{itoa} with sign & Phase 2 (\code{kprintf}) \\
\code{bitmap.c} & \code{bitmap_set}, \code{bitmap_clear}, \code{bitmap_test}, \code{bitmap_find_first_clear} & Phase 5 (frame allocator) \\
\code{ringbuf.c} & fixed-size byte ring buffer: \code{rb_put}, \code{rb_get}, \code{rb_empty}, \code{rb_full} & Phase 4 (keyboard) \\
\code{list.h} & intrusive doubly linked list: \code{list_add}, \code{list_del}, \code{list_for_each} & Phases 7--8 (heap, tasks) \\
\bottomrule
\end{tabularx}
\end{center}

\begin{lst}{sh}{terminal}
gcc -std=gnu11 -Wall -Wextra -g -fsanitize=address,undefined \
    -Iinclude tests/test_lib.c lib/*.c -o build/test_lib && ./build/test_lib
\end{lst}
\verify{All asserts pass with sanitizers enabled, including edge cases: zero-length copies, overlapping
\code{memmove}, \code{utoa(0)}, \code{itoa(INT_MIN)}, a full ring buffer, the last bit of a bitmap.}

\task{Read your first assembly}
You will learn to write assembly in Phase 1, but you can start reading it today. Write three tiny functions (add two
ints, sum an array, call another function) and look at what the compiler makes of them. Compare \code{-O0} with
\code{-O2}, and keep the output in your journal as a phrase book.

\begin{lst}{sh}{terminal}
i686-elf-gcc -std=gnu11 -ffreestanding -O0 -S -masm=intel add.c -o add.s  # read add.s
i686-elf-gcc -std=gnu11 -ffreestanding -O0 -c add.c -o add.o
i686-elf-objdump -d -M intel add.o     # machine code next to the instructions
i686-elf-readelf -S add.o              # the sections: .text .data .bss ...
i686-elf-nm add.o                      # the symbols the linker will see
\end{lst}
\verify{You can point to where each function starts, where its argument is read from the stack, and which register
holds the return value.}

\task{Smoke-test QEMU and GDB}
Start the emulated PC with nothing to boot: \code{qemu-system-i386}. You should see the SeaBIOS firmware try every
device and report that there is no bootable disk. That screen is the \emph{exact} point where your code will take
over in Phase 1. Also open the QEMU monitor (\code{Ctrl+Alt+2} in the window, or start with \code{-monitor stdio}) and
run \code{info registers}.
\verify{You can find the value of \code{EIP} and \code{CR0} in the \code{info registers} output.}

\task{Start the lab journal}
Create \code{docs/journal.md}. For every session, write down the date, what you tried, what happened, and what you
learned. When a bug takes more than an hour, write down the hypothesis before each attempt. This feels slow and
saves weeks.

\section{Testing and debugging}

Most problems in this phase are environment problems. The table covers the common ones.

\begin{pitfalls}
\code{nasm: command not found} in VS Code, but it works in a normal terminal & Flatpak sandbox (Task~0.2) & Use the host terminal profile or a normal terminal. \\
GCC build stops with missing \code{gmp.h} or \code{mpfr.h} & Build dependencies missing & Install \code{gmp-devel mpfr-devel libmpc-devel}, rerun \code{configure}. \\
\code{configure: error: building in source directory} & Built inside the source tree & Use a separate \code{build-gcc} directory. \\
\code{i686-elf-gcc} found in one terminal, not another & \code{PATH} set only in that shell & Put the export in \code{~/.bashrc}; open a new terminal. \\
Sanitizer reports \code{heap-buffer-overflow} in your \code{strncpy} & Off-by-one on the terminator & Good: this is why you test on Linux first. Fix it now, not inside a kernel. \\
\end{pitfalls}

\section{Milestone}

\begin{milestonebox}[Workshop ready]
\begin{checklist}
  \item \code{i686-elf-gcc}, \code{i686-elf-ld}, \code{nasm}, \code{qemu-system-i386} and \code{gdb} run from your editor's terminal.
  \item The Turk-OS repository exists with the skeleton folders and at least one commit.
  \item The warm-up library passes its tests under AddressSanitizer and UBSan.
  \item You have compiled C to assembly and can explain the output of a small function.
  \item You have seen the SeaBIOS ``no bootable device'' screen and used \code{info registers}.
  \item \code{docs/journal.md} has its first entries.
\end{checklist}
\end{milestonebox}

\begin{questionbox}
\begin{enumerate}
  \item What are the three big ideas of \OSTEP, and which Turk-OS phases correspond to each?
  \item What is the difference between kernel mode and user mode, and who enforces it?
  \item Which C headers are still available in freestanding mode, and why exactly those?
  \item What can the compiler do to a non-\code{volatile} pointer access that would break a device driver?
  \item Why does the \code{.bss} section take no space in the ELF file? Who zeroes it?
  \item Round \code{0x00123456} up to the next 4\,KiB boundary by hand.
  \item Name two things the host GCC assumes that are false inside a kernel.
\end{enumerate}
\end{questionbox}

\section{Going further}

If you finish early, set up \code{clangd} or the VS Code C/C++ extension with a \code{compile_commands.json} so
the editor understands your cross-compiled code (\code{bear -- make} generates the file). Work through the Makefile
chapter of the GNU Make manual, especially automatic variables (\code{$@}, \code{$<}, \code{$^}) and pattern rules.
Then read \MOS\,\S1.4 \emph{The Operating System Zoo} for the wider landscape.

\nextphase{1}{x86 Assembly and the Boot Process}
\booklegend
\end{document}
